\uuid{6Vmh}
\exo7id{471}
\titre{exo7 471}
\auteur{cousquer}
\organisation{exo7}
\datecreate{2003-10-01}
\isIndication{false}
\isCorrection{false}
\chapitre{Propriétés de R}
\sousChapitre{Maximum, minimum, borne supérieure}
\cpge{mpsi-mp2i.nombres-reels-suites.b, pcsi.nombres-reels-suites.a, ptsi.nombres-reels-suites.a, mpsi-mp2i.nombres-reels-suites.d, pcsi.nombres-reels-suites.c, ptsi.nombres-reels-suites.c, bcpst1.nombres-reels, bcpst1.suites-reelles}
\module{Analyse}
\niveau{L1}
\difficulte{}

\contenu{
\question{
Soit $E$ l'ensemble des réels de la forme $\frac{n-1/n}{n + 1/n}$ avec
$n \in\mathbb{N}^*$. L'ensemble $E$ admet-il une borne inférieure, une borne
supérieure, un plus grand élément, un plus petit élément ?
}
}
